\section{Locked model and observables}\label{sec:model}

Let \(e_0,\ldots,e_5\) be the standard basis of \(\C^6\), with source \(e_0\)
and target \(e_5\).  The unperturbed real-symmetric Hamiltonian is
\begin{equation}\label{eq:H0}
H_0=\begin{pmatrix}
0&1&1&1&1&0\\
1&0&1&1&0&1\\
1&1&0&0&1&-1\\
1&1&0&0&1&1\\
1&0&1&1&0&-1\\
0&1&-1&1&-1&0
\end{pmatrix}.
\end{equation}
For \(d=(d_1,d_2,d_3,d_4)\), set
\begin{equation}
D(d)=\diag(0,d_1,d_2,d_3,d_4,0),\qquad H(d)=H_0+D(d).
\end{equation}
The fixed-ray specialisation is \(d=\e v\), where \(v\in\R^4\) is selected
independently of \(\e\).  A moving direction \(v=v(\eta)\) is a different,
two-parameter object; it is not smuggled into a fixed-ray theorem.

Three observables must be distinguished:
\begin{align}
G(z;d)&=e_5^{\mathsf T}(zI-H(d))^{-1}e_0,\label{eq:resolvent}\\
M_k(d)&=e_5^{\mathsf T}H(d)^ke_0,\label{eq:moment}\\
K(d,t)&=e_5^{\mathsf T}e^{-itH(d)}e_0
       =\sum_{k\ge0}\frac{(-i)^k}{k!}M_k(d)t^k.\label{eq:amplitude}
\end{align}
The numerator of \(G\), the moment series and the amplitude are linked but not
identical.  In particular, the factor \((-i)^k/k!\) can change a face
cancellation equation when different time depths share a pulled-back weight.
The passage from the finite moment coefficients to this actual matrix
exponential is itself part of the formal development: M8 proves the convergent
identity \cref{eq:leanbridge} for every fixed complex ray, perturbation
parameter and time.

Write
\begin{equation}\label{eq:threeforms}
L=d_1-d_2+d_3-d_4,\qquad Q=d_2d_4-d_1d_3,
\end{equation}
and
\begin{equation}\label{eq:R}
R=d_1d_2(d_3-d_4)+d_3d_4(d_1-d_2).
\end{equation}
On a ray \(d=\e(a,b,c,d)\) we write
\(\ell=a-b+c-d\) and retain \(Q=bd-ac\); the letter \(R\) then denotes the
cubic coefficient in \cref{eq:R}, not a remainder, a Rees algebra, a ray
parameter, or the cubic moment.  This notation lock prevents the principal
symbol collision in the source record.

\begin{figure}[ht]
\centering\begin{tikzpicture}[scale=.95,>=Latex,
  v/.style={circle,draw,thick,fill=white,minimum size=8mm,font=\small},
  p/.style={v,fill=blue!10}, posedge/.style={thick}, neg/.style={thick,dashed,red!70!black}]
\node[v] (s) at (-3,0) {$0$};
\node[p] (a) at (-1.1,1.5) {$1\;[d_1]$};
\node[p] (b) at (-1.1,-1.5) {$2\;[d_2]$};
\node[p] (c) at (1.1,1.5) {$3\;[d_3]$};
\node[p] (d) at (1.1,-1.5) {$4\;[d_4]$};
\node[v] (t) at (3,0) {$5$};
\foreach \x in {a,b,c,d}\draw[posedge] (s)--(\x);
\draw[posedge] (a)--(b); \draw[posedge] (a)--(c); \draw[posedge] (a)--(t);
\draw[posedge] (b)--(d); \draw[neg] (b)--(t);
\draw[posedge] (c)--(d); \draw[posedge] (c)--(t);
\draw[neg] (d)--(t);
\node[font=\scriptsize,align=left] at (0,-2.25) {solid: $+1$\qquad dashed red: $-1$};
\end{tikzpicture}

\caption{The locked six-state channel.  Solid lines carry (+1), dashed red
lines carry (-1), and the four internal vertices carry diagonal coordinates.}
\label{fig:model}
\end{figure}

\begin{proposition}[Baseline darkness]\label{prop:baseline-dark}
For \(d=0\), \(K(0,t)=0\) for every \(t\in\R\), equivalently
\(M_k(0)=0\) for every \(k\ge0\).
\end{proposition}
\begin{proof}
Let $S$ be either of the two signed involutions defined by the model.  It
satisfies $S^2=I$, $SH_0=H_0S$, $Se_0=e_0$ and
$e_5^{\mathsf T}S=-e_5^{\mathsf T}$.  Hence, for every $k$,
\[
M_k(0)=e_5^{\mathsf T}H_0^ke_0
=e_5^{\mathsf T}S H_0^k S e_0=-M_k(0),
\]
so $M_k(0)=0$ over characteristic zero.  Termwise substitution in the
globally convergent exponential series gives $K(0,t)=0$.
\end{proof}

\noindent\textbf{Formal status.}
The matrix identities, abstract signed-certificate implication and all-depth
conclusion are \evidence{KERNEL-CHECKED} as
\path{M3A.rOrig_involutive}, \path{M3A.rOrig_commutes_H0},
\path{M3A.selectedMoment_eq_zero_of_signedCertificate} and
\path{M3A.baseline_moments_dark}.  M8 checks the matrix-exponential passage;
no finite moment cutoff is used as a proof.

\begin{remark}
The proposition is about the channel, not about the persistence of any
particular certificate.  A perturbation may destroy a signed involution while
the transfer remains exactly zero for another reason.
\end{remark}
